\section{Introduction}
Let $n$ be a positive integer and $Fl_n$ the full-flag variety of ${\mathbb C}^n$. As a set, $Fl_n$ is the collection of nested sequences of linear subspaces of ${\mathbb C}^n$ given as follows:
\begin{align*}
 Fl_n = \{V_{\bullet}=(V_1\subset V_2\subset \cdots \subset V_{n}={\mathbb C}^n) \mid \dim_{{\mathbb C}}V_i=i \ (1\le i\le n)\}.
\end{align*}
Let $N$ be an $n\times n$ regular nilpotent matrix viewed as a linear map $N\colon {\mathbb C}^n\rightarrow {\mathbb C}^n$. The Peterson variety $Pet_{n}\subseteq Fl_n$ is defined by
\begin{align*}
 Pet_{n} \coloneqq \{V_{\bullet} \in Fl_n \mid NV_i\subseteq V_{i+1} \ (1\le i\le n-1)\},
\end{align*}
where $NV_i$ denotes the image of $V_i$ under the map $N\colon {\mathbb C}^n\rightarrow {\mathbb C}^n$. It was introduced by Dale Peterson to study the quantum cohomology ring of $Fl_n$, and it has since appeared in several contexts (e.g.\ \cite{AHMMS,Ba,ha-ho-ma,Kostant,Rietsch06}). 

For a permutation $w\in\mathfrak{S}_n$, let $X_w\subseteq Fl_n$ be the Schubert variety associated with~$w$, and $\Omega_w\subseteq Fl_n$ the dual Schubert variety associated with $w$. We denote by~$[n-1]$ the set of integers $1,2,\ldots,n-1$. For a subset $J\subseteq[n-1]$, let $w_J\in\mathfrak{S}_n$ be the longest element of the Young subgroup $\mathfrak{S}_J$ of the permutation group $\mathfrak{S}_n$ associated with $J$ (see Section~\ref{subsect: combi} for details), and set
\begin{align*}
X_{J} 
\coloneqq X_{w_J}\cap Pet_{n} \quad \text{and} \quad \Omega_{J} 
\coloneqq \Omega_{w_J}\cap Pet_{n}.
\end{align*}
Then $X_J$ and $\Omega_J$ in $Pet_{n}$ play roles analogous to those of Schubert varieties and dual Schubert varieties in $Fl_n$, and provide important information about the topology of~$Pet_{n}$.

In this paper, we construct an additive basis $\{\varpi_J\mid J\subseteq[n-1]\}$ of the integral cohomology ring $H^*(Pet_{n};{\mathbb Z})$ which reflects the geometry of $X_J$ and $\Omega_J$ (Theorem~\ref{thm: Z-basis of cohomology of Peterson}). As a consequence, we may consider the structure constants for the multiplication rule:
\begin{align}\label{intro:100}
 \varpi_J\cdot \varpi_K = \sum_{L\subseteq[n-1]} d_{JK}^L \varpi_L , \qquad d_{JK}^L\in {\mathbb Z}.
\end{align}
It turns out that all $d_{JK}^L$ are non-negative integers, and we give a geometric proof of this positivity by using that of nef line bundles over $Pet_{n}$ (Proposition~\ref{prop: positivity}). We also provide a manifestly positive combinatorial formula for $d_{JK}^L$ (Theorem~\ref{theorem:LeftRight_diagram}) in terms of \textit{left-right diagrams}, which we introduce in this paper. 

To find our formula for $d_{JK}^L$, 
we prove several properties of the cohomology classes $\varpi_J$ which are inherited from the geometry of $X_J$ and $\Omega_J$.
In particular, writing $\Omega_J$ as an intersection of divisors on $Pet_{n}$ provides the geometric idea behind our formula for $d_{JK}^L$ in terms of left-right diagrams.

We also show that our basis $\{\varpi_J\mid J\subseteq[n-1]\}$ coincides with the additive basis of the cohomology ring $H^*(Pet_{n};{\mathbb C})$ with ${\mathbb C}$-coefficients introduced by Harada--Tymoczko~\cite{HaTy}. Their basis is obtained by taking restriction of certain Schubert classes to $Pet_{n}$, and it is called the \textit{Peterson Schubert basis}. It has been studied by Bayegan--Harada~\cite{BaHa}, Drellich~\cite{Dre2} and Goldin--Gorbutt~\cite{GoGo}. In \cite{GoGo}, Goldin and Gorbutt gave combinatorial formulas for the structure constants of Harada--Tymoczko's basis (in a certain equivariant setting) which are manifestly positive and integral. 
Thus, after taking the non-equivariant limit, their formulas and ours both describe the same structure constants, but these formulas have different perspectives; 
their approach is mostly combinatorial whereas our approach is based on the geometry of $X_J$ and $\Omega_J$. We include a short comparison of their formulas and ours in Section~\ref{sec: relations to other works}.

Interestingly, our computations match with those of Berget--Spink--Tseng~\cite[Sect.~7]{Berget-Spink-Tseng} on a certain subring of the cohomology ring of a toric variety which is called the permutohedral variety.
One of their results can be interpreted as a formula describing the structure constants $d_{JK}^L$ as products of \textit{mixed Eulerian numbers} which were introduced and studied by Postnikov~\cite{Postnikov09}.
With this connection in mind, our formula (Theorem~\ref{theorem:LeftRight_diagram}) for $d_{JK}^L$ can also be thought as computing some products of mixed Eulerian numbers by using the geometry of $Pet_{n}$.
We explain this connection in Section~\ref{sec: relations to other works} including the relations with the works of Nadeau--Tewari~\cite{Nadeau-Tewari} and the second author~\cite{Horiguchi21}.

\longthanks{We are grateful to Mikiya Masuda for his support and encouragement. We also thank the anonymous referees for reading this paper very carefully. This research is supported in part by Osaka City University Advanced Mathematical Institute (MEXT Joint Usage/Research Center on Mathematics and Theoretical Physics): the topology and combinatorics of Hessenberg varieties. The first author is supported in part by JSPS Grant-in-Aid for Early-Career Scientists: 18K13413. The second author is supported in part by JSPS Grant-in-Aid for Young Scientists: 19K14508. The fourth author is supported in part by NSFC: 11901218.}

\section{Basic notations}\label{sec: Basic notations}
In this section, we recall some terminology which will be used in this paper.

\subsection{Combinatorics on the Dynkin diagram of type A}\label{subsect: combi}
Let $n(\ge2)$ be a positive integer. We use the notation $[n-1]\coloneqq\{1,2,\ldots,n-1\}$, and we regard it as the set of vertices of the Dynkin diagram of type $A_{n-1}$ for the rest of the paper. Namely, two vertices $i,j\in[n-1]$ are connected by an edge if and only if $|i-j|=1$. See Figure~\ref{pic: Dynkin diagram}.

\begin{figure}[htbp]
\[\includegraphics[width=5.3cm]{figures/pic1}\]\vspace{-20pt}
\caption{The Dynkin diagram of type $A_{n-1}$.}
\label{pic: Dynkin diagram}
\end{figure}

We regard each subset $J\subseteq[n-1]$ as a full-subgraph of the Dynkin diagram appearing above. We may decompose it into the connected components: 
\begin{align}\label{eq: connected components}
 J = J_1 \sqcup J_2 \sqcup \cdots \sqcup J_{m},
\end{align}
where each $J_k\ (1\le k\le m)$ is the set of vertices of a maximal connected subgraph of~$J$. To determine each $J_k$ uniquely, we require that the maximal element of~$J_k$ is less than the minimal element of $J_{k'}$ when $k< k'$.

\begin{exam}
Let $n=10$ and $J=\{1,2,4,5,6,9\}$. Then we have 
\begin{align*}
 J = \{1,2\} \sqcup\{4,5,6\} \sqcup\{9\} = J_1\sqcup J_2\sqcup J_3.
\end{align*}
\end{exam}

For $J\subseteq[n-1]$, one has the Young subgroup given by
\begin{align*}
 \mathfrak{S}_J \coloneqq \mathfrak{S}_{J_1}\times \mathfrak{S}_{J_2}\times \cdots \mathfrak{S}_{J_m}
 \subseteq \mathfrak{S}_n,
\end{align*}
where each $\mathfrak{S}_{J_k}$ is the subgroup of $\mathfrak{S}_n$ generated by the simple reflections $s_i$ for all~$i\in J_k$. Let $w_J$ be the longest element of $\mathfrak{S}_J$, i.e.
\begin{align}\label{eq: def of wJ}
 w_J \coloneqq w_0^{(J_1)} w_0^{(J_2)} \cdots w_0^{(J_{m})} \in \mathfrak{S}_J, 
\end{align}
where each $w_0^{(J_k)}$ is the longest element of the permutation group $\mathfrak{S}_{J_k}$ $(1\le k\le m)$. 

\begin{exam}\label{eg: wJ}
{\rm
Let $n=10$ and $J=\{1,2\} \sqcup\{4,5,6\} \sqcup\{9\}$ as above. By identifying the permutation $w_J$ with its permutation matrix, we have 
\begin{align*}
 w_J 
 = w_0^{(J_1)} w_0^{(J_2)} w_0^{(J_3)} 
 = (s_1s_2s_1) (s_4s_5s_6s_4s_5s_4) (s_9)
 =
 \left(
 \begin{array}{@{\,}ccc|cccc|c|cc@{\,}}\vspace{-3pt}
     & & 1 & & & & & & & \\ 
     & 1 & & & & & & & & \vspace{-3pt}\\ 
    1 & & & & & & & & & \\ \hline
     & & & & & & 1& & &  \vspace{-3pt}\\
     & & & & & 1 & & & & \vspace{-3pt}\\
     & & & & 1 & & & & & \vspace{-3pt}\\
     & & & 1 & & & & & & \\ \hline 
     & & & & & & & 1 & & \vspace{-2pt}\\ \hline
     & & & & & & & & & 1 \vspace{-3pt}\\
     & & & & & & & & 1 & 
 \end{array}
 \right).
\end{align*}}
\end{exam}

We can identify each permutation $w\in\mathfrak{S}_n$ with its permutation flag $V_{\bullet}\in Fl_n$ defined by $V_i =\langle e_{w(1)}, e_{w(2)},\ldots, e_{w(i)}\rangle$ for $1\le i\le n$, where $e_1,e_2,\ldots,e_n$ denotes the standard basis of ${\mathbb C}^n$, and the right hand side is the linear subspace of ${\mathbb C}^n$ spanned by $e_{w(1)}, e_{w(2)},\ldots, e_{w(i)}$. Using this identification, we explain how the permutations $w_J$ are related to the Peterson variety. Let ${\rm GL}_n({\mathbb C})$ be the general linear group of invertible $n\times n$ complex matrices. Let $T\subseteq {\rm GL}_n({\mathbb C})$ be the maximal torus consisting of diagonal matrices. Let us identify ${\mathbb C}^{\times}$ with a subgroup of $T$ as follows: 
\begin{align}\label{eq: C-starsubgroup}
{\mathbb C}^{\times} = 
\left\{
\left.
\begin{pmatrix}
g & & & \\
 & g^2 & & \\
 & & \ddots & \\
 & & & g^n 
\end{pmatrix}
\in T\ 
\right| \ g\in {\mathbb C}^{\times}
\right\}.
\end{align}
The flag variety $Fl_n$ admits a natural ${\rm GL}_n({\mathbb C})$-action by regarding each element $g\in {\rm GL}_n({\mathbb C})$ as an automorphism $g\colon {\mathbb C}^n\rightarrow {\mathbb C}^n$. Restricting this ${\rm GL}_n({\mathbb C})$-action on $Fl_n$ to the above subgroup ${\mathbb C}^{\times}$, it is well-known that the fixed point set $(Fl_n)^{{\mathbb C}^{\times}}$ is the set of the permutation flags, i.e. $(Fl_n)^{{\mathbb C}^{\times}}=\mathfrak{S}_n$ (e.g.~\cite[proof of Lemma~2 in Sect.~10.1]{fult97}). It is straightforward to see that this ${\mathbb C}^{\times}$-action on $Fl_n$ preserves $Pet_{n}$, and it was shown in \cite{HaTy} that the fixed point set $(Pet_{n})^{{\mathbb C}^{\times}}$ is given by
\begin{align}\label{eq: fixed points of Pet}
(Pet_{n})^{{\mathbb C}^{\times}} =\{w_J \in\mathfrak{S}_n \mid J\subseteq [n-1]\}.
\end{align}
Because of this relation, the combinatorics of $w_J$ will be important to understand the structure of $Pet_{n}$.

For $1\le i\le n-1$, let $s_i\in\mathfrak{S}_n$ be the simple reflection which interchanges $i$ and~$i+1$. We denote by $\le$ the Bruhat order on $\mathfrak{S}_n$, that is, we have $u\le v$ $(u,v\in\mathfrak{S}_n)$ if and only if a reduced expression of $u$ is a subword of a reduced expression of $v$.
\begin{lemma}\label{lem: So-san's lemma 1}
For $J\subseteq[n-1]$ and $1\le i\le n-1$, we have 
$s_i\le w_J$ if and only if $i\in J$.
\end{lemma}

\begin{proof}
Recall that $w_J$ is the product of longest elements in $\mathfrak{S}_{J_k}$ for $1\le k\le m$:
\begin{align*}
 w_J = w_0^{(J_1)} w_0^{(J_2)} \cdots w_0^{(J_{m})}.
\end{align*}
Since each $w_0^{(J_k)}$ ($1\le k\le m$) preserves the decomposition \eqref{eq: connected components}, it follows that the length of $w_J$ is the same as the sum of the lengths of $w_0^{(J_k)}$ for $1\le k\le m$. Thus, the products of reduced expressions of $w_0^{(J_k)}$ for $1\le k\le m$ give a reduced expression of~$w_J$. Here, an arbitrary reduced expression of $w_0^{(J_k)}$ contains a simple reflection $s_i$ if and only if $i\in J_k$. Therefore, it follows that a reduced expression of $w_J$ contains $s_i$ if and only if $i\in J$ (see Example~\ref{eg: wJ}). This implies the desired claim. 
\end{proof}

The following claim appears in \cite[Lemma 6]{Insko-Tymoczko}, but we give a proof for the reader.
\begin{lemma}\label{lem: So-san's lemma}
For $J,J'\subseteq[n-1]$, we have
\begin{align}\label{eq: subavr 15}
w_{J'} \le w_{J} \quad \text{if and only if} \quad J'\subseteq J .
\end{align}
\end{lemma}

\begin{proof}
We first prove that $J'\subseteq J$ under the assumption $w_{J'} \le w_{J}$. For this, take an arbitrary element $i\in J'$. By the previous lemma, we have $s_i\le w_{J'}$. Combining this with the assumption $w_{J'} \le w_{J}$, we obtain that $s_i\le w_{J}$. Thus, it follows that $i\in J$ by the previous lemma again. 

We next prove that $w_{J'} \le w_{J}$ under the assumption $J'\subseteq J$. Take the decomposition 
\begin{align*}
 J' = J'_1 \sqcup J'_2 \sqcup \cdots \sqcup J'_{m'}
\end{align*}
into the connected components as in \eqref{eq: connected components}. Since each $J'_{\ell}\ (1\le \ell \le m')$ is connected, it is contained in some connected component $J_k$ of $J$. This leads us to consider a map
\begin{align*}
 \varphi \colon \{1,2,\ldots,m'\} \rightarrow \{1,2,\ldots,m\}
\end{align*}
which we define by the conditions $J'_i\subseteq J_{\varphi(i)}$ for $1\le i\le m'$. Then we have that
\begin{align*}
 \bigsqcup_{\varphi(i)=k} J'_i \subseteq J_{k}
\end{align*}
by the definition of the map $\varphi$. This implies that 
\begin{align*}
 \prod_{\varphi(i)=k} w_0^{(J'_i)} \le w_0^{(J_k)}
 \qquad \text{in $\mathfrak{S}_{J_k}$}
\end{align*}
since $w_0^{(J_k)}$ is the longest permutation in $\mathfrak{S}_{J_k}$. Recalling that each $J_k$ is a connected component of $J$, these inequalities for $1\le k\le m$ imply that $w_{J'}\le w_{J}$.
\end{proof}

\begin{exam}
\emph{
Let $n=8$, $J'=\{1,4,5,7\}$ and $J=\{1,2,4,5,6,7\}$ so that we have $J'\subseteq J$. In this case, we have
\begin{align*}
 J'=\{1\}\sqcup\{4,5\}\sqcup\{7\}=J'_1\sqcup J'_2\sqcup J'_3
 \quad \text{and} \quad
 J=\{1,2\}\sqcup\{4,5,6,7\}=J_1\sqcup J_2,
\end{align*}
and hence
\begin{align*}
 w_{J'}= (s_1)(s_4s_5s_4)(s_7) \le (s_1s_2s_1)(s_4s_5s_6s_7s_4s_5s_6s_4s_5s_4)=w_J.
\end{align*}
}
\end{exam}

\subsection{Hessenberg varieties}\label{subsect: Hessenberg}
In this subsection, we briefly recall the notion of Hessenberg varieties in type A. They are (possibly reducible) subvarieties of the flag variety $Fl_n$, and will appear in the next section. A function $h\colon [n] \to [n]$ is a \textit{Hessenberg function} if it satisfies the following two conditions:
\begin{enumerate}
\item[(i)] $h(1) \leq h(2) \leq \cdots \leq h(n)$,
\item[(ii)] $h(j) \geq j$ for all $j \in [n]$. 
\end{enumerate}
Note that $h(n)=n$ by definition. We may identify a Hessenberg function $h$ with a configuration of (shaded) boxes on a square grid of size $n \times n$ which consists of boxes in the $i$-th row and the $j$-th column satisfying $i \leq h(j)$ for $i, j\in[n]$, as we illustrate in the following example.

\begin{exam}\label{example:HessenbergFunction}
\emph{
Let $n=5$. The Hessenberg function $h\colon[5]\rightarrow [5]$ given by 
\[(h(1),h(2),h(3),h(4),h(5))=(2,3,3,5,5)\]
corresponds to the configuration of the shaded boxes drawn in Figure $\ref{pic:stair-shape}$.}
\begin{figure}[htbp]
\[\includegraphics[width=2.8cm]{figures/pic2}\]
\vspace{-20pt}
\caption{The configuration of the shaded boxes corresponding to $h$.}
\label{pic:stair-shape}
\end{figure}
\end{exam}

For an $n\times n$ matrix $X$ considered as a linear map $X\colon{\mathbb C}^n \to {\mathbb C}^n$ and a Hessenberg function $h\colon [n] \to [n]$, the \textit{Hessenberg variety} associated with $X$ and $h$ (\cite{DeMari-Procesi-Shayman,ty}) is defined as
\begin{equation} \label{eq:DefinitionHessenbergVariety}
\text{Hess}(X,h)=\{V_\bullet \in Fl_n \mid XV_j \subseteq V_{h(j)} \ \mbox{for all} \ j\in[n] \}.
\end{equation}
We note that $\text{Hess}(X,h)\cong \text{Hess}(gXg^{-1},h)$ for all $g\in{\rm GL}_n({\mathbb C})$ so that we may always assume that the matrix $X$ is in Jordan canonical form.
Let $N$ be the $n\times n$ regular nilpotent matrix in Jordan canonical form, i.e. 
\begin{align}\label{eq: def of N}
N =
\begin{pmatrix} 
0 & 1 & & &\\ 
& 0 & 1 & &\\
&   & \ddots & \ddots &\\
&  & & 0 & 1\\
&  & & & 0
\end{pmatrix}.
\end{align}
Then $\text{Hess}(N,h)$ is called a \textit{regular nilpotent Hessenberg variety}. It is known that $\text{Hess}(N,h)$ is an irreducible projective variety of (complex) dimension $\sum_{i=1}^n (h(i)-i)$ (\cite[Lemma 7.1]{an-ty} and \cite{ty}). When $h(i)=i+1$ for all $1\le i\le n-1$, it is clear that $\text{Hess}(N,h)=Pet_{n}$ by definition. 
In particular, we obtain the well-known formula 
\[\dim_{{\mathbb C}}Pet_{n}=n-1.\]
For Hessenberg functions $h,h'\colon[n]\rightarrow [n]$, it is clear that if $h(i)\le h'(i)$ for $1\le i\le n$ then $\text{Hess}(N,h)\subseteq \text{Hess}(N,h')$. For example, the Hessenberg function $h\colon[5]\rightarrow[5]$ given in Example~\ref{example:HessenbergFunction} defines a $3$-dimensional regular nilpotent Hessenberg variety $\text{Hess}(N,h)$ which is contained in $Pet_{5}(\subseteq Fl_5)$. 

\section{Geometric constructions}\label{sec: geometric constructions}
In this section, we introduce two kinds of subvarieties $X_J$ and $\Omega_J$ in $Pet_{n}$ for each $J\subseteq[n-1]$, and we establish geometric properties of them. They will play important roles in constructing an additive basis of the integral cohomology ring $H^*(Pet_{n};{\mathbb Z})$ in the next section.

\subsection{Analogue of Schubert varieties in the Peterson variety}
For $w\in \mathfrak{S}_n$, 
let $X_{w}\subseteq Fl_n$ be the Schubert variety associated with $w$ and $\Omega_w\subseteq Fl_n$ the dual Schubert variety associated with $w$ (\cite[Sect.~10]{fult97}).
We note that $\dim_{{\mathbb C}} X_{w} =\codim_{{\mathbb C}} (\Omega_w,Fl_n)= \ell(w)$ and $X_w\cap \Omega_w=\{w\}$, where $\ell(w)$ is the length of $w$.
\begin{defi}
For $J\subseteq [n-1]$, we define
\begin{equation}
\begin{split}\label{eq: subavr 10}
X_{J} 
\coloneqq X_{w_J}\cap Pet_{n} \quad \text{and} \quad \Omega_{J} 
\coloneqq \Omega_{w_J}\cap Pet_{n},
\end{split}
\end{equation}
where $w_J\in\mathfrak{S}_n$ is the permutation defined in $\eqref{eq: def of wJ}$.
\end{defi}

Peterson~\cite{Peterson} studied a particular open affine subset of $\Omega_J$ to construct the quantum cohomology ring of $Fl_n$ (c.f.~\cite{Kostant,Rietsch06}). Also, Insko~\cite{Insko} and Insko--Tymoczko~\cite{Insko-Tymoczko} studied $X_J$ to show the injectivity of the homomorphism $H_*(Pet_{n};{\mathbb Z})\rightarrow H_*(Fl_n;{\mathbb Z})$. It turns out that $X_J$ and $\Omega_J$ in $Pet_{n}$ play roles analogous to those of Schubert varieties and dual Schubert varieties in $Fl_n$. As an illustrating property, we begin with the following claim. Recall that we have $X_{w}\cap \Omega_{v}\ne\emptyset$ in $Fl_n$ if and only if $w\ge v$. 

\begin{prop}\label{prop: point intersection 2}
For $J,J'\subseteq[n-1]$, we have
\begin{align*}
X_{J}\cap \Omega_{J'} \neq \emptyset
\quad \text{if and only if} \quad
J \supseteq J'.
\end{align*}
Moreover, when $J=J'$, we have $X_J \cap \Omega_J = \{w_{J}\}$.
\end{prop}


\begin{proof}
If $X_{J}\cap \Omega_{J'} \neq \emptyset$, then we have $(X_{w_J}\cap\Omega_{w_{J'}})\cap Pet_{n}\ne\emptyset$ by definition. Note that $(X_{w_J}\cap\Omega_{w_{J'}})\cap Pet_{n}$ is complete, and it is preserved by the ${\mathbb C}^{\times}$-action on $Pet_{n}$ described in Section~\ref{sec: Basic notations}. Thus, it follows that it contains a ${\mathbb C}^{\times}$-fixed point (e.g.~\cite[Chap.~VIII, Sect.~21.2]{Humphreys}). Since we have
\begin{align*}
((X_{w_J}\cap\Omega_{w_{J'}})\cap Pet_{n})^{{\mathbb C}^{\times}} 
= (X_{w_J}\cap\Omega_{w_{J'}})^{{\mathbb C}^{\times}} \cap (Pet_{n})^{{\mathbb C}^{\times}},
\end{align*}
we see that there exists a ${\mathbb C}^{\times}$-fixed point $w_K\in Pet_{n}$ (see \eqref{eq: fixed points of Pet}) such that $w_K\in X_{w_J}$ and $w_K\in\Omega_{w_{J'}}$. The former condition implies that $w_J\ge w_K$, and the latter condition implies that $w_K \ge w_{J'}$ (\cite[Sect.~10.2 and~10.5]{fult97}). Thus, we obtain $w_J\ge w_{J'}$, and it follows that $J\supseteq J'$ from Lemma~\ref{lem: So-san's lemma}.

If $J \supseteq J'$, then we have $w_J\ge w_{J'}$ by Lemma~\ref{lem: So-san's lemma}. This implies that $w_{J}\in X_{w_J}\cap\Omega_{w_{J'}}$, and hence $X_{J}\cap \Omega_{J'}\ne\emptyset$ follows.
\end{proof}

A distinguished property of $X_J$ is that it is a regular nilpotent Hessenberg variety for a certain Hessenberg function. Let us explain this in the following. For each $J\subseteq[n-1]$, there is a natural Hessenberg function which is determined by $J$ as follows. Let $h_J\colon[n]\rightarrow[n]$ be a function given by
\begin{align}\label{eq: def of hJ}
 h_J (i) =
 \begin{cases}
  i+1 \quad &\text{if $i\in J$}\\
  i &\text{if $i\notin J$}
 \end{cases}
\end{align}
for $1\le i\le n$. Then $h_J$ is a Hessenberg function, and we have $\text{Hess}(N,h_J)\subseteq Pet_{n}$ since the Hessenberg function for $Pet_{n}$ is given by $h(i)=i+1$ for $1\le i\le n-1$ as we saw in Section~\ref{subsect: Hessenberg}.

\begin{exam}\label{ex: hJ}
{\rm 
Let $n=10$ and $J= \{1,2\} \sqcup\{4,5,6\} \sqcup\{9\} = J_1\sqcup J_2\sqcup J_3$. Then the configuration of boxes of $h_J$ is given in Figure $\ref{pic: example of hJ}$. Compare the figure with the permutation matrix of $w_J$ in Example~$\ref{eg: wJ}$.}
\begin{figure}[htbp]
\[\includegraphics[width=4.3cm]{figures/pic3}\]
\caption{The Hessenberg function $h_J$.}
\label{pic: example of hJ}
\end{figure}
\end{exam}

As we mentioned above, Insko--Tymoczko~\cite{Insko-Tymoczko} studied $X_J$, and they proved most of the following claim. Recall that $X_J$ is defined to be the intersection $X_{w_J}\cap Pet_{n}$ where $X_{w_J}$ is the Schubert variety associated with $w_J$. 

\begin{prop}\label{prop: XJ and Hess}
For $J\subseteq [n-1]$, we have
\begin{align}\label{eq: two equalities}
 X_J = \overline{X^{\circ}_{w_J}\cap Pet_{n}} = \text{Hess}(N,h_J)
\end{align}
where $N$ is the regular nilpotent matrix given in \eqref{eq: def of N}, and $X^{\circ}_{w_J}$ is the Schubert cell associated with $w_J$. In particular, we have $\dim_{{\mathbb C}} X_J = |J|$.
\end{prop}

To prove this, we need the following lemma. Let $X^{\circ}_{w}\subseteq Fl_n$ be the Schubert cell associated with $w$ and and $\Omega^{\circ}_w\subseteq Fl_n$ the dual Schubert cell associated with $w$.

\begin{lemma}\label{lem: intersection with dual cell}
The following are equivalent.
\begin{enumerate}
 \item $X^{\circ}_w\cap Pet_{n}\ne\emptyset$
 \item $\Omega^{\circ}_w\cap Pet_{n}\ne\emptyset$
 \item $w\in Pet_{n}$ $($i.e. $w=w_J$ for some $J\subseteq[n-1])$
\end{enumerate}
\end{lemma}

\begin{proof}
It is clear that (3) implies (1). To see that (1) implies (3), take an  element $z\in X^{\circ}_w\cap Pet_{n}\ne\emptyset$. Since $X^{\circ}_w\cap Pet_{n}\subseteq X^{\circ}_w$ is preserved under the ${\mathbb C}^{\times}$-action on $X^{\circ}_w$, it follows that $t\cdot z\in X^{\circ}_w\cap Pet_{n}$ for all $t\in {\mathbb C}^{\times}$. Noticing that $X^{\circ}_w\cap Pet_{n}\subseteq X^{\circ}_w$ is a closed subset, we have 
\begin{align}\label{eq: limit}
\lim_{t\rightarrow0}t\cdot z\in X^{\circ}_w\cap Pet_{n}.
\end{align}
Under the standard identification $X^{\circ}_w = {\mathbb C}^{\ell(w)}$ (c.f.~\cite[Sect.\ 10.2]{fult97}), the ${\mathbb C}^{\times}$-action on $X^{\circ}_w$ is identified with a linear action with positive weights. Thus it follows that $\lim_{t\rightarrow0}t\cdot z=0$, which corresponds to $w\in X^{\circ}_w$ (c.f.~\cite[proof of Lemma~5]{Insko-Tymoczko}). Thus it follows that $w\in Pet_{n}$ by \eqref{eq: limit}. The equivalence of (2) and (3) follows by an argument similar to that for the equivalence of (1) and (3).
\end{proof}

\begin{proof}[Proof of Proposition~$\ref{prop: XJ and Hess}$]
Let us first prove that 
\begin{align}\label{eq: first inclusion}
 X^{\circ}_{w_J}\cap Pet_{n} \subseteq \text{Hess}(N,h_J)
 \qquad \text{for each $J\subseteq [n-1]$}.
\end{align}
For this, take an arbitrary element $V_{\bullet}\in X^{\circ}_{w_J}\cap Pet_{n}$. Then we have
\begin{align*}
 NV_i \subseteq V_{i+1} \qquad (1\le i\le n-1).
\end{align*}
To see that $V_{\bullet}\in \text{Hess}(N,h_J)$, we need to show that 
\begin{align}\label{eq: NVp}
 NV_p \subseteq V_{p} \qquad \text{for }p\notin J.
\end{align}
Since we are assuming $V_{\bullet}\in X^{\circ}_{w_J}\cap Pet_{n}$, the flag $V_{\bullet}$ lies in the Schubert cell $X^{\circ}_{w_J}$. Here, the permutation $w_J$ is a product of the longest permutations of the symmetric group of smaller ranks as given in \eqref{eq: def of wJ}. Thus it follows (from e.g.~\cite[Sect.~10.2]{fult97}) that 
\begin{align*}
 V_p 
 = 
 \langle e_{1}, e_{2}, \ldots, e_{p}\rangle
 \qquad \text{for }p\notin J ,
\end{align*}
where $e_1,e_2,\ldots,e_n$ is the standard basis of ${\mathbb C}^n$. Since $Ne_1 =0$ and $Ne_i =e_{i-1}$ for~$2\le i\le n$, it is clear that \eqref{eq: NVp} follows. Thus we obtain \eqref{eq: first inclusion}.

Now, we prove the claim \eqref{eq: two equalities} of this proposition. Since we have $X_{w_J}=\bigsqcup_{v\le w_J}  X^{\circ}_{v}$, it follows that
\begin{align}\label{eq: decomp of XJ}
 X_{J}
 = X_{w_J}\cap Pet_{n}
 = \bigsqcup_{v\le w_J} ( X^{\circ}_{v}\cap Pet_{n} )
 = \bigsqcup_{J'\subseteq J} ( X^{\circ}_{w_{J'}}\cap Pet_{n} ), 
\end{align}
where the last equality follows from Lemmas \ref{lem: So-san's lemma} and \ref{lem: intersection with dual cell}. For each intersection $X^{\circ}_{w_{J'}}\cap Pet_{n}$ in the right-most side, we have that $X^{\circ}_{w_{J'}}\cap Pet_{n}\subseteq \text{Hess}(N,h_{J'})$ by \eqref{eq: first inclusion}. 
The condition $J'\subseteq J$ implies that $h_{J'}(i)\le h_J(i)$ for $1\le i\le n$, and hence we have $\text{Hess}(N,h_{J'})\subseteq \text{Hess}(N,h_{J})$. Combining this with the previous inclusion, we see that
\begin{align*}
 X^{\circ}_{w_{J'}}\cap Pet_{n}\subseteq \text{Hess}(N,h_J) 
\end{align*}
in \eqref{eq: decomp of XJ}. Thus it follows that
\begin{align}\label{eq: XJ in HesshJ}
 X_{J} \subseteq \text{Hess}(N,h_J).
\end{align}
Note that $X^{\circ}_{w_J}\cap Pet_{n} \subseteq X_{J}(=X_{w_J}\cap Pet_{n})$ by definition, and hence we have that $\overline{X^{\circ}_{w_J}\cap Pet_{n}} \subseteq X_{J}$ by taking the closure. Combining this with \eqref{eq: XJ in HesshJ}, we obtain that
\begin{align}\label{eq: two inclusions}
 \overline{X^{\circ}_{w_J}\cap Pet_{n}} \subseteq X_{J} \subseteq \text{Hess}(N,h_J).
\end{align}
In this sequence, both sides have the same dimension.  This is because we have $\dim_{{\mathbb C}} \overline{X^{\circ}_{w_J}\cap Pet_{n}} = |J|$ from \cite[Lemma~9]{Insko-Tymoczko} and $\dim_{{\mathbb C}} \text{Hess}(N,h_J)=|J|$ from \cite[Lemma~7.1]{an-ty}. Since $\text{Hess}(N,h_J)$ is irreducible, the two inclusions in \eqref{eq: two inclusions} are equalities. This completes the proof.
\end{proof}

Combining Proposition~\ref{prop: XJ and Hess} and a result of Drellich \cite[Theorem 4.5]{Dre1}, we may express $X_J$ as a product of Peterson varieties of smaller ranks as follows. 
Let $J = J_1 \sqcup J_2 \sqcup \cdots \sqcup J_{m}$ be the decomposition into the connected components. Then, by definition, we have $w_J=w_0^{(J_1)} w_0^{(J_2)} \cdots w_0^{(J_{m})}$, where each $w_0^{(J_k)}$ is a product of longest elements in $\mathfrak{S}_{J_k}$ for $1\le k\le m$. Hence it follows that the Schubert variety $X_{w_J}\subseteq Fl_n$ associated with $w_J$ is isomorphic to the product of the flag varieties of smaller ranks:
\begin{align*}
 X_{w_J} = \prod_{k=1}^{m} X_{w_0^{(J_k)}} \cong \prod_{k=1}^{m} Fl_{n_k}, 
\end{align*}
where we set
\begin{align*}
 n_k\coloneqq |J_k|+1 \qquad \text{for $1\le k\le m$}.
\end{align*}
By restricting this isomorphism to $X_J = X_{w_J}\cap Pet_{n}$, it follows from Proposition~\ref{prop: XJ and Hess} and \cite[Theorem 4.5]{Dre1} that $X_{J}$ is isomorphic to a product of Peterson varieties of smaller ranks. 
\begin{coro}\label{lem: decomp of PetJ}
For $J\subseteq[n-1]$, we have
\begin{align*}
 X_{J} = \prod_{k=1}^{m} X_{J_k} \cong
 \prod_{k=1}^{m} Pet_{n_k},
\end{align*}
where $J = J_1 \sqcup J_2 \sqcup \cdots \sqcup J_{m}$ is the decomposition into the connected components and $n_k= |J_k|+1$ $(1\le k\le m)$.
\end{coro}

\begin{exam}\label{ex: product decomp}
{\rm 
Let $n=10$ and $J =\{1,2\} \sqcup\{4,5,6\} \sqcup\{9\} = J_1\sqcup J_2\sqcup J_3$. The representation matrix of $w_J$ is given in Example~$\ref{eg: wJ}$, and we have 
\begin{align*}
&X_{w_J} \cong Fl_3\times Fl_4\times Fl_2, \\
&X_{J} \cong Pet_{3}\times Pet_{4}\times Pet_{2}.
\end{align*}
}
\end{exam}

Compared to Schubert varieties and dual Schubert varieties in $Fl_n$, the structures of $X_J$ and $\Omega_J$ in $Pet_{n}$ are rather simple as we explain below.
To begin with, we make the following definition.
\begin{defi}
For $1\le i\le n-1$, let
\begin{equation}
\begin{split}\label{eq: div 10}
D_{i} \coloneqq X_{[n-1]\setminus\{i\}}
\quad \text{and} \quad
E_{i} \coloneqq \Omega_{\{i\}}.
\end{split}
\end{equation}
where $X_{[n-1]\setminus\{i\}}$ and $\Omega_{\{i\}}$ are defined in \eqref{eq: subavr 10}.
\end{defi}

\begin{lemma}\label{lem: divisors}
For $1\le i\le n-1$, the following hold. 
\begin{enumerate}
 \item $D_{i}$ and $E_{i}$ have codimension $1$ in $Pet_{n}$.
 \item $D_{i}\cap E_{i} = \emptyset$.
\end{enumerate}
\end{lemma}

\begin{proof}
For (1), we have $\dim_{{\mathbb C}}D_i=\dim_{{\mathbb C}} X_{[n-1]\setminus\{i\}} =n-2$ by Proposition~\ref{prop: XJ and Hess}. We also have 
\begin{align*}
E_{i} = \Omega_{\{i\}} = \Omega_{w_{\{i\}}}\cap Pet_{n}= \Omega_{s_i}\cap Pet_{n} .
\end{align*}
It is well-known that $\Omega_{s_i}$ is irreducible and it has complex codimension $1$ in $Fl_n$ (\cite[Sect.\ 10.2]{fult97}). Hence $\Omega_{s_i}$ in $Fl_n$ is locally cut out by a single function. We also know that $\Omega_{s_i}\cap Pet_{n}$ is a non-empty proper subset of $Pet_{n}$ since we have $s_i\in \Omega_{s_i}\cap Pet_{n}$ and $\text{id}=w_{\emptyset}\in Pet_{n}\setminus\Omega_{s_i}$. Thus, it follows that $\dim_{{\mathbb C}}E_{i} =n-2$. For (2), the claim follows from Proposition~\ref{prop: point intersection 2}.
\end{proof}

In the next subsection, we will see that $D_i$ and $E_i$ are divisors\footnote{In this paper, a divisor on a variety $Y$ always means the support of an effective Cartier divisor on $Y$, i.e. the zero locus of a section of a line bundle over $Y$.} on $Pet_{n}$. The following claim means that $X_J$ and $\Omega_J$ can be described as intersections of divisors on $Pet_{n}$.

\begin{prop}\label{prop: divisor expression}
For $J\subseteq [n-1]$, we have
\begin{align}\label{eq: div exp 10}
X_{J} =\bigcap_{i\notin J} D_i
\quad \text{and} \quad
\Omega_{J} = \bigcap_{i\in J} E_i .
\end{align}
\end{prop}

\begin{proof}
By \eqref{eq: decomp of XJ}, we have
\begin{align*}
 X_{w_{[n-1]\setminus\{i\}}}\cap Pet_{n} = \bigsqcup_{J'\subseteq [n-1]\setminus\{i\}} (X^{\circ}_{w_{J'}}\cap Pet_{n}).
\end{align*}
This implies from the definition of $D_i$ that
\begin{align*}
 \bigcap_{i\notin J} D_i
 = \bigcap_{i\notin J} (X_{w_{[n-1]\setminus\{i\}}}\cap Pet_{n}) 
 = \bigsqcup_{J'\subseteq J} (X^{\circ}_{w_{J'}}\cap Pet_{n}).
\end{align*}
Combining this with \eqref{eq: decomp of XJ}, we obtain the desired claim for $X_J$. An argument similar to this proves that 
\begin{align*}
 \bigcap_{i\in J} E_i
 = \bigcap_{i\in J} (\Omega_{w_{\{i\}}}\cap Pet_{n}) 
 = \bigsqcup_{J''\supseteq J} (\Omega^{\circ}_{w_{J''}}\cap Pet_{n})
 = \Omega_{w_J}\cap Pet_{n}
 = \Omega_J
\end{align*}
by Lemmas \ref{lem: So-san's lemma} and \ref{lem: intersection with dual cell}.
\end{proof}

\begin{exam}\label{eg: intersection}
{\rm 
Let $n=9$ and $J=\{2,3,4\}\sqcup\{7,8\}$ so that $[n-1]\setminus J=\{1\}\sqcup\{5,6\}$. Then we have 
\begin{align*}
X_{J} = D_1\cap D_5\cap D_6
\qquad \text{and} \qquad
\Omega_{J} = E_2\cap E_3\cap E_4\cap E_7\cap E_8.
\end{align*}
}
\end{exam}

\subsection{Defining equations of \texorpdfstring{$X_J$}{XJ} and \texorpdfstring{$\Omega_J$}{OmegaJ}}\label{subsec: equation of X and Omega}
Let $B\subseteq {\rm GL}_n({\mathbb C})$ be the Borel subgroup of upper-triangular matrices in ${\rm GL}_n({\mathbb C})$. Then we have the standard identification $Fl_n={\rm GL}_n({\mathbb C})/B$. For a complex $B$-representation space $V$, we have the associated complex vector bundle\footnote{We take the $B$-action on the product ${\rm GL}_n({\mathbb C})\times V$ so that $[g,v]=[gb,b^{-1}\cdot v]$ in the quotient.} over ${\rm GL}_n({\mathbb C})/B$:
\begin{align*}
{\rm GL}_n({\mathbb C})\times^B V \rightarrow {\rm GL}_n({\mathbb C})/B
\quad ; \quad
[g,v]\mapsto gB.
\end{align*}
For a weight $\mu\colon T\rightarrow {\mathbb C}^{\times}$, we obtain $\mu\colon B\rightarrow {\mathbb C}^{\times}$ by composing that with the canonical projection $B\twoheadrightarrow T$, which we also denote by the same symbol. We denote by ${\mathbb C}_{\mu}={\mathbb C}$ the corresponding 1-dimensional representation space of $B$. Set
\begin{align}\label{eq: def of Lmu 2}
L_{\mu} = {\rm GL}_n({\mathbb C})\times^B {\mathbb C}_{\mu}^* = {\rm GL}_n({\mathbb C})\times^B {\mathbb C}_{-\mu},
\end{align}
where ${\mathbb C}_{\mu}^*$ is the dual representation space of ${\mathbb C}_{\mu}$. We also denote the restriction $L_{\mu}|_{Pet_{n}}$ by the same symbol $L_{\mu}$ when there is no confusion.

Let us introduce two representations of $B$ associated with each $J\subseteq [n-1]$ as follows. For $1\le i\le n-1$, let $\varpi_i\colon T\rightarrow {\mathbb C}^{\times}$ be the $i$-th fundamental weight of $T$ given by $\text{diag}(t_1,t_2,\ldots,t_n)\mapsto t_1t_2\cdots t_i$. For $J\subseteq [n-1]$, we obtain a representation space of $T$ given by a direct sum
\begin{align*}
\bigoplus_{i\in J} {\mathbb C}_{\varpi_i}^*.
\end{align*}
Through the canonical projection $B\twoheadrightarrow T$, we regard this as a representation of $B$. To introduce the other representation of $B$ associated with $J$, 
let $\alpha_i$ $(1\le i\le n-1)$ be the $i$-th simple root defined as a weight $\alpha_i\colon T\rightarrow {\mathbb C}^{\times}$ given by $\text{diag}(t_1,t_2,\ldots,t_n)\mapsto t_{i}t_{i+1}^{-1}$. Let $H_J\subseteq \mathfrak{gl}_n({\mathbb C})$ be the Hessenberg subspace (c.f. \cite[Sect.~2]{ty}) corresponding to the Hessenberg function $h_J$ defined in \eqref{eq: def of hJ}, that is,
\begin{align}\label{eq: HJ}
&H_{J} \coloneqq \mathfrak{b}\oplus\bigoplus_{i\in J} \mathfrak{g}_{-\alpha_i} \subseteq \mathfrak{gl}_n({\mathbb C}),
\end{align}
where $\mathfrak{b}=\text{Lie}(B)$ is the Lie algebra of $B$ and each $\mathfrak{g}_{-\alpha_{i}}$ is the standard root space of~$\mathfrak{gl}_n({\mathbb C})$ associated with the $i$-th negative simple root $-\alpha_i$ $(1\le i\le n-1)$. Since $H_J$ is preserved by the adjoint action of $B$ on $\mathfrak{gl}_n({\mathbb C})$, the quotient space
\begin{align*}
H_{[n-1]}/H_J
\end{align*}
is a representation space of $B$. Now, these two representations of $B$ induce the following vector bundles over $Fl_n$: 
\begin{align*}
&U_{J} 
\coloneqq {\rm GL}_n({\mathbb C})\times^B \left(H_{[n-1]}/H_J\right), \\
&V_{J} 
\coloneqq {\rm GL}_n({\mathbb C})\times^B \left(\bigoplus_{i\in J} {\mathbb C}_{\varpi_i}^*\right).
\end{align*}
If there is no confusion, we denote the restrictions of $U_{J}$ and $V_{J}$ on $Pet_{n}$ by the same symbol. Note that we have
\begin{equation}
\begin{split}\label{eq: rank of U and J}
&\rank U_{J}=(n-1)-|J|, \\
&\rank V_{J}=|J|.
\end{split}
\end{equation}

Recall that 
\begin{align*}
Pet_{n} = \{gB \in {\rm GL}_n({\mathbb C})/B \mid g^{-1}Ng\in H_{[n-1]}\}
\end{align*}
(c.f.~\cite{Insko-Tymoczko,Rietsch06} or~\cite[Sect.~2]{ty}).
Thus, the following map gives a section of $U_{J}$ over $Pet_{n}$:
\begin{align*}
\phi_J \colon  Pet_{n} \rightarrow U_{J}
\quad ; \quad 
\phi_J(gB) = [g, [g^{-1}Ng]],
\end{align*}
where $[g^{-1}Ng]\in H_{[n-1]}/H_J$ is the class represented by $g^{-1}Ng \in H_{[n-1]}$. For $1\le i\le n-1$, let 
\begin{align*}
\text{det}_i \colon {\rm GL}_n({\mathbb C})\rightarrow {\mathbb C}_{\varpi_i}(={\mathbb C})
\end{align*} 
be the function which takes the leading principal minor of order $i$. This is a $B$-equivariant function with respect to the multiplication of $B$ on ${\rm GL}_n({\mathbb C})$ from the right. Thus, we have a well-defined section 
\begin{align*}
\psi_J \colon  Pet_{n} \rightarrow V_J
\quad ; \quad
gB\mapsto \Big[g,\sum_{i\in J}\text{det}_i(g)\Big].
\end{align*}
The following claim means that $X_J$ and $\Omega_J$ in $Pet_{n}$ are defined by the equation $\phi_J=0$ and $\psi_J=0$, respectively.

\begin{prop}\label{prop: zero locus}
For $J\subseteq[n-1]$, we have 
\begin{align*}
 X_J=Z(\phi_J) \quad \text{and} \quad \Omega_J =Z(\psi_J) ,
\end{align*}
where $Z(\phi_J)$ and $Z(\psi_J)$ denote the zero loci of the sections $\phi_J$ and $\psi_J$, respectively.
\end{prop}

\begin{proof}
Since the defining condition of $\text{Hess}(N,h_{J})$ is precisely that $g^{-1}Ng\in H_{J}$ (e.g.~\cite[Sect.~2]{ty}), it is clear that we have $X_J=Z(\phi_J)$. It is known that 
\begin{align*}
\Omega_{s_i} = \{gB \in {\rm GL}_n({\mathbb C})/B \mid \text{det}_i(g)=0\}
\end{align*}
as subsets of $Fl_n$ (c.f.~\cite[Proposition~9 in Sect.~10.6]{fult97}). Thus, it follows that
\begin{align*}
Z(\psi_{\{i\}}) = \Omega_{s_i}\cap Pet_{n}=E_i.
\end{align*}
Now, we obtain that
\begin{align*}
Z(\psi_J) = \bigcap_{i\in J} Z(\psi_{\{i\}}) = \bigcap_{i\in J} E_i = \Omega_{J}
\end{align*}
by Proposition~\ref{prop: divisor expression}.
\end{proof}

\begin{coro}\label{cor: cartier}
For $1\le i\le n-1$, $D_i$ and $E_i$ are divisors on $Pet_{n}$.
\end{coro}

\begin{proof}
By \eqref{eq: rank of U and J}, it follows that $U_{[n-1]\setminus\{i\}}$ is a line bundle, and we have $D_i=X_{[n-1]\setminus\{i\}}=Z(\phi_{[n-1]\setminus\{i\}})$ by the previous proposition. This means that $D_i$ is a divisor on $Pet_{n}$. Similarly, $E_i=\Omega_{\{i\}}=Z(\psi_{\{i\}})$ is a divisor on $Pet_{n}$ since $V_{\{i\}}$ is a line bundle by \eqref{eq: rank of U and J}.
\end{proof}

\begin{coro}\label{cor: codim of OmegaJ}
For $J\subseteq[n-1]$, we have $\codim_{{\mathbb C}} \Omega_J=|J|$ in $Pet_{n}$.
\end{coro}

\begin{proof}
Recall that $\Omega_J =Z(\psi_J)$ and $\rank V_{J}=|J|$. This means that $\Omega_J$ in $Pet_{n}$ is locally cut out by $|J|$ functions, and it follows that each irreducible component of $\Omega_J$ has codimension at most $|J|$ in $Pet_{n}$ (\cite[Proposition 14.1]{Fulton}). This implies that
\begin{align}\label{eq: ineq from equations}
\dim_{{\mathbb C}}\Omega_J\ge (n-1)-|J|
\end{align}
since $\dim_{{\mathbb C}}Pet_{n}=n-1$ and $\Omega_J=\Omega_{w_J}\cap Pet_{n}\ne\emptyset$. We show that the equality holds by induction on $(n-1)-|J|$. When $J=[n-1]$, we have $\Omega_{[n-1]}=\Omega_{w_0}\cap Pet_{n}=\{w_0\}$ so that the claim is obvious. Let $J\subsetneq[n-1]$, and assume by induction that $\dim_{{\mathbb C}}\Omega_K= (n-1)-|K|$ for $J\subsetneq K\subseteq[n-1]$. We prove that $\dim_{{\mathbb C}}\Omega_J= (n-1)-|J|$. Take an element $i\in[n-1]\setminus J$, and set $K=J\sqcup\{i\}$. Then we have 
\begin{align}\label{eq: decreasing sequence}
\Omega_{K}=\Omega_J\cap E_i \subseteq \Omega_J
\end{align}
by Proposition~\ref{prop: divisor expression}. This means that $\Omega_{K}$ is the zero locus of the section $\psi_{\{i\}}$ of the line bundle $V_{\{i\}}$ restricted over $\Omega_{J}$. Thus we see that 
\begin{align}\label{eq: zero-locus ineq}
\dim_{{\mathbb C}}\Omega_{K} \ge \dim_{{\mathbb C}} \Omega_{J} -1,
\end{align}
which follows by applying \cite[Proposition 14.1]{Fulton} to each irreducible component of $\Omega_{J}$. Namely, the dimension decreases at most by 1 in \eqref{eq: decreasing sequence}. Since we have $\dim_{{\mathbb C}}\Omega_K= (n-1)-|K|$ by the inductive hypothesis, we can rewrite \eqref{eq: zero-locus ineq} as
\begin{align*}
\dim_{{\mathbb C}}\Omega_{J} \le \dim_{{\mathbb C}}\Omega_K +1 = (n-1) - |K| + 1 = (n-1) - |J|.
\end{align*}
Combining this with \eqref{eq: ineq from equations}, we obtain the desired equality.
\end{proof}

\begin{rema}\label{rem: decomposition into line bundles}
The vector bundles $U_J$ and $V_J$ are decomposed into line bundles as
\begin{align*}
U_J = \bigoplus_{i\notin J} L_{\alpha_i}
\qquad \text{and}\qquad
V_J = \bigoplus_{i\in J} L_{\varpi_i}.
\end{align*} 
The first equality follows since we have $H_{[n-1]}/H_J\cong \oplus_{i\notin J}{\mathbb C}_{-\alpha_i}$ as representations of~$B$, where the right hand side is a direct sum representation. These decompositions can be viewed as the analogue of Proposition $\ref{prop: divisor expression}$ in the language of vector bundles over $Pet_{n}$.
\end{rema}

\section{The cohomology ring of \texorpdfstring{$Pet_{n}$}{Petn}}\label{sec: cohomology}
In this section, we construct an additive basis of the integral cohomology ring $H^*(Pet_{n};{\mathbb Z})$ by incorporating the geometry established in the previous section.
We also introduce the structure constants of the basis, and provide a geometric proof for their positivity.

For an algebraic variety $X$ which admits a paving by complex affine spaces (\cite[Definition~2.1]{ty}),
an irreducible Zariski closed subset $Y$ of $X$ has its fundamental cycle (as a reduced scheme) in $H_{2d}(Y;{\mathbb Z})$, where $d=\dim_{{\mathbb C}}Y$.
By abusing notation, we use the same symbol for its image in $H_{2d}(X;{\mathbb Z})$ under the induced map $i_*\colon H_{2d}(Y;{\mathbb Z})\rightarrow H_{2d}(X;{\mathbb Z})$, where $i$ is the inclusion map $i\colon Z\hookrightarrow X$. See \cite[Appendix~B]{fult97} or \cite[Chap.~19]{Fulton} for the details.

\subsection{The homology group of \texorpdfstring{$Pet_{n}$}{Petn}}
Recall that we have a decomposition of $Fl_n$ by Schubert cells:
\begin{align*}
  Fl_n = \bigsqcup_{w\in\mathfrak{S}_n} X^{\circ}_{w}.
\end{align*}
This induces a set-theoretic decomposition
\begin{align*}
 Pet_{n} = \bigsqcup_{J\subseteq[n-1]} (X^{\circ}_{w_J}\cap Pet_{n}),
\end{align*}
by Lemmas~\ref{lem: So-san's lemma} and \ref{lem: intersection with dual cell}. It is known from \cite{Peterson,ty} that each $X^{\circ}_{w_J}\cap Pet_{n}$ is isomorphic to an affine cell ${\mathbb C}^{|J|}$ and that these cells form a paving by affines. Recall also from Proposition~\ref{prop: XJ and Hess} that we have
\begin{align*}
 X_J = \overline{X^{\circ}_{w_J}\cap Pet_{n}}
\end{align*}
and that $\dim_{{\mathbb C}}X_J=|J|$ for $J\subseteq[n-1]$. This implies that the cycles represented by~$X_J$ for $J\subseteq[n-1]$ form a ${\mathbb Z}$-basis of the homology group $H_*(Pet_{n};{\mathbb Z})$. 

\begin{prop}\label{prop: homology basis}
$($\cite{Peterson,ty}$)$
For each $J\subseteq [n-1]$, we have $[X_J]\in H_{2|J|}(Pet_{n};{\mathbb Z})$, and the set
$\{ [X_J] \mid J\subseteq [n-1]\}$ is a ${\mathbb Z}$-basis of $H_*(Pet_{n};{\mathbb Z})$; 
\begin{align*}
 H_*(Pet_{n};{\mathbb Z}) = \bigoplus_{J\subseteq [n-1]} {\mathbb Z} [X_J] .
\end{align*}
\end{prop}

\begin{exam}\label{ex: homology basis}
{\rm 
Let $n=4$ so that $[n-1]=\{1,2,3\}$. Then we have 
\begin{align*}
H_*(Pet_{n};{\mathbb Z}) &= {\mathbb Z} [X_{\emptyset}]\oplus({\mathbb Z} [X_{\{1\}}]\oplus{\mathbb Z} [X_{\{2\}}]\oplus{\mathbb Z} [X_{\{3\}}])\\
&\hspace{75pt}\oplus({\mathbb Z} [X_{\{1,2\}}]\oplus{\mathbb Z} [X_{\{1,3\}}]\oplus{\mathbb Z} [X_{\{2,3\}}])\oplus{\mathbb Z} [X_{\{1,2,3\}}].
\end{align*}}
\end{exam}

\begin{rema}
Compared to $X^{\circ}_{w_J}\cap Pet_{n}(\cong {\mathbb C}^{|J|})$, the geometry of $\Omega^{\circ}_{w_J}\cap Pet_{n}$ is known to encode the quantum cohomology ring of a partial flag variety specified by~$J$ (see~\cite{Peterson, Rietsch06}).
\end{rema}

\subsection{The cohomology group of \texorpdfstring{$Pet_{n}$}{Petn}}
For each weight $\mu\colon T\rightarrow {\mathbb C}^{\times}$, we constructed a line bundle $L_{\mu}$ over $Fl_n$ in Section~\ref{subsec: equation of X and Omega}.
Recall also that $\alpha_i$ and $\varpi_i$ are the $i$-th simple root and the $i$-th fundamental weight of $T$ $(1\le i\le n-1)$, respectively.
It is well-known that we have an isomorphism
\begin{align*}
\bigoplus_{i=1}^{n-1} {\mathbb Z}\varpi_i
\stackrel{\cong}{\rightarrow} H^2(Fl_n;{\mathbb Z}) 
\quad ; \quad
\mu \mapsto e(L_{\mu}),
\end{align*}
where we regard each $\mu=a_1\varpi_1+\cdots+a_{n-1}\varpi_{n-1}\ (a_1,\ldots,a_{n-1}\in{\mathbb Z})$ as the weight $\mu\colon T\rightarrow {\mathbb C}^{\times}$ given by $\text{diag}(t_1,\ldots,t_n)\mapsto t_1^{a_1}(t_1t_2)^{a_2}\cdots(t_1\cdots t_{n-1})^{a_{n-1}}$.
Let~$i\colon Pet_{n}\hookrightarrow Fl_n$ be the inclusion map. Insko~\cite{Insko} proved that the induced homomorphism $i_*\colon H_*(Pet_{n};{\mathbb Z})\rightarrow H_*(Fl_n;{\mathbb Z})$ is an injection whose image is a direct summand of $H_*(Fl_n;{\mathbb Z})$. This means that the map $i_*\colon H_2(Pet_{n};{\mathbb Z}) \rightarrow H_2(Fl_n;{\mathbb Z})$ on degree $2$ is an isomorphism since $\rank H_2(Fl_n;{\mathbb Z})=\rank H_2(Pet_{n};{\mathbb Z}) =n-1$ (see Proposition~\ref{prop: homology basis}), and hence the restriction map 
\begin{align*}
i^*\colon H^2(Fl_n;{\mathbb Z})  \stackrel{\cong}{\rightarrow} H^2(Pet_{n};{\mathbb Z}) 
\end{align*}
on degree $2$ cohomology group is an isomorphism. By combining these isomorphisms, we obtain that $\bigoplus_{i=1}^{n-1} {\mathbb Z}\varpi_i\cong H^2(Pet_{n};{\mathbb Z})$. In the rest of the paper, we identify $\bigoplus_{i=1}^{n-1} {\mathbb Z}\varpi_i$ with $H^2(Pet_{n};{\mathbb Z})$ through this isomorphism, and we use the same symbol $\mu\in H^2(Pet_{n};{\mathbb Z})$ for the element  $e(L_{\mu})$ by abusing notation. For example, we write
\begin{align*}
\alpha_i = e(L_{\alpha_i})
\qquad \text{and} \qquad
\varpi_i = e(L_{\varpi_i})
\end{align*}
as elements in $ H^2(Pet_{n};{\mathbb Z})$.

In Section~\ref{subsec: equation of X and Omega}, we constructed vector bundles $U_J$ and $V_{J}$ over $Pet_{n}$. Adopting the above notation, we may express the Euler class $e(V_J)$ as a monomial of $\varpi_i$ for each $i\in J$. Namely, for $J\subseteq[n-1]$, we have
\begin{align}\label{eq: monomial exp of e(VJ)}
e(V_J) = \prod_{i\in J} \varpi_i
\end{align}
since the vector bundle $V_{J}$ decomposes into line bundles as follows:
\begin{align*}
V_{J} 
= {\rm GL}_n({\mathbb C})\times^B \left(\bigoplus_{i\in J} {\mathbb C}_{\varpi_i}^*\right)
= \bigoplus_{i\in J} L_{\varpi_i}.
\end{align*}

For $J\subseteq[n-1]$, take the decomposition $J = J_1 \sqcup J_2 \sqcup \cdots \sqcup J_{m}$ into the connected components. We set 
\begin{align}\label{eq: def of mJ}
m_J \coloneqq |J_1| ! |J_2| ! \cdots |J_m| !.
\end{align}

\begin{defi} \label{defi:varpi}
For $J\subseteq[n-1]$, let
\begin{align*}
 \varpi_J \coloneqq \frac{1}{m_J} 
 e(V_J)= \frac{1}{m_J} \prod_{i\in J} \varpi_i ,
\end{align*}
where $m_J$ is defined in \eqref{eq: def of mJ}.
\end{defi}

The cohomology class $\varpi_J$ is defined to be an element of $H^{2|J|}(Pet_{n};{\mathbb Q})$, but we will show that it belongs to the integral cohomology group $H^{2|J|}(Pet_{n};{\mathbb Z})$.

\begin{exam}
{\rm 
Let $n=9$ and $J=\{2,3,4,7,8\}$ so that $J=\{2,3,4\}\sqcup\{7,8\}$ is the decomposition into the connected components. Then we have 
\begin{align*}
 \varpi_J 
 = \frac{1}{3!2!}(\varpi_2\varpi_3\varpi_4)(\varpi_7\varpi_8)
 =\frac{1}{12}\varpi_2\varpi_3\varpi_4\varpi_7\varpi_8 .
\end{align*}
Compare this with Example~$\ref{eg: intersection}$.}
\end{exam}

\begin{rema}
We also have 
\begin{align*}
e(U_J) = \prod_{i\notin J} \alpha_i 
\end{align*}
$($c.f. Remark~$\ref{rem: decomposition into line bundles}$$)$. These decompositions of $e(U_J)$ and $e(V_J)$ can be viewed as the cohomological analogue of Proposition $\ref{prop: divisor expression}$. 
\end{rema}

The main purpose of this subsection is to prove that the set of cohomology classes $\{\varpi_J \mid J\subseteq[n-1]\}$ forms a module basis of the integral cohomology group $H^*(Pet_{n};{\mathbb Z})$. We will state this in Theorem~\ref{thm: Z-basis of cohomology of Peterson}, and we devote the rest of this subsection for its proof.

\begin{lemma}\label{lem: vanishing product}
For $1\le i\le n-1$, we have 
\begin{align*}
\alpha_i \varpi_i = 0 \quad \text{in $H^4(Pet_{n},{\mathbb Z})$}.
\end{align*}
\end{lemma}

\begin{proof}
Notice that $\alpha_i \varpi_i$ is the Euler class of the rank 2 vector bundle $U_{[n-1]\setminus\{i\}}\oplus V_{\{i\}}=L_{\alpha_i}\oplus L_{\varpi_i}$ (c.f. Remark~\ref{rem: decomposition into line bundles}). From Section~\ref{subsec: equation of X and Omega}, we have the section $\phi_{[n-1]\setminus\{i\}}+\psi_{\{i\}}$ of this bundle whose zero locus is $Z(\phi_{[n-1]\setminus\{i\}})\cap Z(\psi_{\{i\}})=D_i\cap E_i$ as we saw in the proof of Corollary~\ref{cor: cartier}.
Now, by Lemma~\ref{lem: divisors}~(2), this is the empty set. Thus,~$\phi_{[n-1]\setminus\{i\}}+\psi_{\{i\}}$ on $Pet_{n}$ is a nowhere-zero section, and hence the Euler class $\alpha_i \varpi_i$ vanishes (see~\cite[Property 9.7]{Milnor-Stasheff}).
\end{proof}

\begin{rema}\label{rem: fundamental rel}
In \cite[Corollary~3.4]{fu-ha-ma} and \cite[Theorem 4.1]{ha-ho-ma}, the equations $\alpha_i \varpi_i = 0$ for $1\le i\le n-1$ appeared as the fundamental relations in the presentation of the cohomology ring $H^*(Pet_{n};{\mathbb C})$. 
\end{rema}

For $1\le i\le n$, let $F_i$ be the $i$-th tautological vector bundle over $Fl_n$ whose fiber at $V_{\bullet}\in Fl_n$ is $V_i$. As a convention, we denote by $F_0$ the trivial sub-bundle of $F_1$ of rank~$0$. The quotient line bundle $L_i\coloneqq F_i/F_{i-1}$ is called the $i$-th tautological line bundle, and we set 
\begin{align} \label{eq:xi}
x_i \coloneqq c_1(L_i^*) \in H^2(Fl_n;{\mathbb Z})
\quad (1\le i\le n),
\end{align}
where we note that $x_1+\cdots+x_n=0$. We will also denote by the same symbol the restriction of $x_i$ to $H^2(Pet_{n};{\mathbb Z})$. It is well-known that for $1\le i\le n-1$, we have 
\begin{equation}\label{eq: pi and x}
\begin{split}
&\alpha_i = x_i-x_{i+1}, \\
&\varpi_i = x_1+x_2+\cdots+x_i.
\end{split}
\end{equation}
For $1\le i<j\le n$, let
\begin{align*}
\alpha_{i,j}\coloneqq \alpha_i+\alpha_{i+1}+\cdots+\alpha_{j-1}=x_i-x_j.
\end{align*}
For a homology cycle $Z\in H_{k}(Fl_n;{\mathbb Z})$ of degree $k$, the Poincar\'{e} dual of $Z$ is the (unique) cohomology class $\gamma\in H^{2d-k}(Fl_n;{\mathbb Z})$ $(d=\dim_{{\mathbb C}}Fl_n)$ satisfying $\gamma\cap[Fl_n]=Z$. In the following lemma, we regard $Fl_{n-1}$ as a subvariety of $Fl_n$ whose flags are contained in the linear subspace of ${\mathbb C}^n$ generated by $e_1,e_2,\ldots,e_{n-1}$.
The claim (2) of the following lemma seems to be well-known, but we provide a proof using Hessenberg varieties for the completeness of the paper.

\begin{lemma}\label{lem: Poincare dual}
The following hold:
\begin{enumerate}
 \item The Poincar\'e dual of $[Pet_{n}]\in H_*(Fl_n;{\mathbb Z})$ is $\prod_{j-i\ge2}\alpha_{i, j}\in H^*(Fl_n;{\mathbb Z})$.
 \item The Poincar\'e dual of $[Fl_{n-1}]\in H_*(Fl_n;{\mathbb Z})$ is $\frac{1}{n}\alpha_{1,n}\alpha_{2,n}\cdots\alpha_{n-1,n}\in H^*(Fl_n;{\mathbb Z})$.
\end{enumerate}
\end{lemma}

\begin{proof}
We first prove the claim (1). Recall from \eqref{eq: def of hJ} and \eqref{eq: HJ} that $H_J\subseteq \mathfrak{gl}_n({\mathbb C})$ for $J\subseteq[n-1]$ is the Hessenberg space corresponding to the Hessenberg function $h_J$. Consider the associated vector bundle 
\begin{align*}
\mathcal{N}\coloneqq {\rm GL}_n({\mathbb C})\times^B (\mathfrak{gl}_n({\mathbb C})/H_{[n-1]})
\end{align*}
over $Fl_n={\rm GL}_n({\mathbb C})/B$. By an argument similar to that used in Section~\ref{subsec: equation of X and Omega}, $Pet_{n}$ can be written as the zero locus of a section of the vector bundle $\mathcal{N}$, and it is shown in \cite[Corollary~3.9]{ab-fu-ze} that the Poincar\'e dual of $[Pet_{n}]$ is the Euler class $e(\mathcal{N})\in H^*(Fl_n;{\mathbb Z})$. It is straightforward to verify that $e(\mathcal{N})=\prod_{j-i\ge2}\alpha_{i, j}$ by the same inductive argument as that in \cite[Sect.\ 4]{ab-fu-ze} using short exact sequences of vector bundles.

Next we prove the claim (2). Let $S$ be an $n\times n$ regular semisimple matrix in diagonal form (i.e. a diagonal matrix with distinct eigenvalues) and $\text{Hess}(S,h_0)$ a regular semisimple Hessenberg variety, where $h_0$ is the Hessenberg function $h_0\colon [n]\rightarrow [n]$ given by
\begin{align*}
h_0(i) \coloneqq
\begin{cases}
n-1 \quad &(1\le i\le n-1) \\
n &(i=n).
\end{cases}
\end{align*}
It is shown in \cite[Sect,\ 3 and Sect.\ 4]{an-ty} that the Poincar\'{e} dual of $[\text{Hess}(S,h_0)]$ is 
\begin{align*}
(x_1-x_n)(x_2-x_n)\cdots (x_{n-1}-x_n)
\in H^*(Fl_n;{\mathbb Z}),
\end{align*}
where the left hand side is equal to $\alpha_{1,n}\alpha_{2,n}\cdots\alpha_{n-1,n}$ by the  definition of $\alpha_{i,j}$. It is also known that $\text{Hess}(S,h_0)$ has $n$ connected components and that all the connected components give the same cycle $[Fl_{n-1}]$ (\cite[Sect.\ 3]{Teff11}). Thus the Poincar\'e dual of~$n[Fl_{n-1}]$ is 
$\alpha_{1,n}\alpha_{2,n}\cdots\alpha_{n-1,n}$, which implies the claim (2).
\end{proof}

\begin{rema}
{\rm
\cite[Corollary~7.2]{an-ty} with the formula for the double Schubert polynomial associated to a dominant permutation given in \cite[p.2613]{an-ty} provides a more general formula than that of Lemma~\ref{lem: Poincare dual}~(1) for regular nilpotent Hessenberg \textit{schemes}, which were not known to be reduced when it was published. After that, \cite{ab-de-ga-ha} proved that they are in fact reduced when they contain $Pet_{n}$ (c.f. \cite[Remark~3.8]{ab-de-ga-ha}), and the formula are now generalized in \cite{ab-fu-ze} for an arbitrary Lie type.
}
\end{rema}

For an (irreducible) projective variety $Y$, we denote the fundamental cycle of $Y$ as $[Y]\in H_{2d}(Y;{\mathbb Z})$, where $d=\dim_{{\mathbb C}} Y$. For a cohomology class $\beta\in H^{2d}(Y;{\mathbb Z})$, we write
\begin{align*}
\int_Y \beta \coloneqq \langle [Y] , \beta \rangle_{Y} \quad (\in{\mathbb Z}), 
\end{align*} 
where the right hand side is the value of the standard paring
\begin{align*}
\langle \ , \ \rangle_{Y} \colon  H_{2d}(Y;{\mathbb Z}) \times H^{2d}(Y;{\mathbb Z}) \rightarrow {\mathbb Z}.
\end{align*}

\begin{prop}\label{lem: Masuda lemma in fund}
We have
\begin{align*}
\int_{Pet_{n}}\varpi_1\varpi_2\cdots \varpi_{n-1} = (n-1)! .
\end{align*}
\end{prop}

\begin{proof}
Let us first prove that
\begin{align}\label{eq: basis in fundamental induction}
\int_{Pet_{n}}\varpi_1\varpi_2\cdots \varpi_{n-1} 
= (n-1) \int_{Pet_{n-1}}\varpi_1\varpi_2\cdots \varpi_{n-2}.
\end{align}
For the $\varpi_{n-1}$ in the left hand side of \eqref{eq: basis in fundamental induction}, notice that
\begin{align*}
\varpi_{n-1} = \frac{1}{n}\sum_{i=1}^{n-1} i\alpha_i ,
\end{align*}
which follows from \eqref{eq: pi and x}. Since we have $\varpi_i\alpha_i=0$ for $1\le i\le n-1$ from Lemma~\ref{lem: vanishing product}, we see that
\begin{align}\label{eq: pi->alpha}
\int_{Pet_{n}}\varpi_1\varpi_2\cdots \varpi_{n-1}
= \frac{n-1}{n}\int_{Pet_{n}}\varpi_1\varpi_2\cdots \varpi_{n-2}\alpha_{n-1}.
\end{align}
By Lemma~\ref{lem: Poincare dual}~(1), the right hand side of \eqref{eq: pi->alpha} can be computed as the following integral over $Fl_n$:
\begin{align*}
\frac{n-1}{n}\int_{Fl_n} \varpi_1\varpi_2\cdots \varpi_{n-2}\alpha_{n-1} \prod_{j-i\ge2}\alpha_{i,j}.
\end{align*}
Since we have $\alpha_{n-1}=\alpha_{n-1,n}$, the last expression can be written as
\begin{align*}
\frac{n-1}{n}\int_{Fl_n} \varpi_1\varpi_2\cdots \varpi_{n-2} (\alpha_{1,n}\alpha_{2,n}\cdots\alpha_{n-1,n})\prod_{\substack{j-i\ge2\\j\ne n}}\alpha_{i, j}.
\end{align*}
By Lemma~\ref{lem: Poincare dual}~(2), we can rewrite this as the following integral over $Fl_{n-1}$:
\begin{align*}
(n-1)\int_{Fl_{n-1}} \varpi_1\varpi_2\cdots \varpi_{n-2} \prod_{\substack{j-i\ge2\\j\ne n}}\alpha_{i, j}.
\end{align*}
Applying Lemma~\ref{lem: Poincare dual}~(1) to $Pet_{n-1}\subseteq Fl_{n-1}$, this can be written as the following integral over $Pet_{n-1}$:
\begin{align*}
(n-1)\int_{Pet_{n-1}} \varpi_1\varpi_2\cdots \varpi_{n-2} .
\end{align*}
Hence we proved \eqref{eq: basis in fundamental induction}.

Now using \eqref{eq: basis in fundamental induction} repeatedly, we obtain that
\begin{align*}
\int_{Pet_{n}}\varpi_1\varpi_2\cdots \varpi_{n-1}
= (n-1)! \int_{Pet_{2}}\varpi_1 .
\end{align*}
Noticing that $Pet_{2}=Fl_2={\mathbb P}^1$, we see that $\varpi_1(=x_1)$ is the first Chern class of the dual of the standard tautological line bundle over ${\mathbb P}^1$ by \eqref{eq:xi}. Thus the integral in the right hand side is equal to $1$, which completes the proof.
\end{proof}

\begin{lemma}\label{lem: geometry of mJ}
For $J\subseteq[n-1]$, we have
\begin{align*}
\int_{X_{J}} e(V_J) = m_J , 
\end{align*}
where $m_J=|J_1| ! |J_2| ! \cdots |J_m| !$ is defined in \eqref{eq: def of mJ}.
\end{lemma}

\begin{proof}
The isomorphism given in Corollary~\ref{lem: decomp of PetJ} induces an isomorphism
\begin{align*}
H^*(X_J;{\mathbb Z}) \stackrel{\cong}{\rightarrow} \bigotimes_{k=1}^{m} H^*(Pet_{n_k};{\mathbb Z})
\end{align*}
which sends
$e(V_J)(=\prod_{i\in J}\varpi_i) \in H^{2|J|}(X_{J};{\mathbb Z})$ to $\bigotimes_{k=1}^{m} \varpi_1\varpi_2\cdots \varpi_{|J_k|}$. It also induces an isomorphism
\begin{align*}
H_*(X_J;{\mathbb Z}) \stackrel{\cong}{\rightarrow} \bigotimes_{k=1}^{m} H_*(Pet_{n_k};{\mathbb Z})
\end{align*}
which sends $[X_{J}]\in H_{2|J|}(X_J;{\mathbb Z})$ to $\bigotimes_{k=1}^{m}[Pet_{n_k}]$. Now the claim follows from Proposition~\ref{lem: Masuda lemma in fund}.
\end{proof}

\begin{prop}\label{prop: duality}
For $J,K\subseteq [n-1]$ such that $|J|=|K|$, the degree of the homology class $[X_{J}]$ is the same as the degree of the Euler class $e(V_K)$, 
and we have
\begin{align*} 
\langle [X_{J}] , e(V_K) \rangle_{Pet_{n}} 
=
\begin{cases}
m_J \quad &\text{if $J=K$}, \\
0 &\text{if $J\ne K$}.
\end{cases}
\end{align*}
\end{prop}

\begin{proof}
Note that we have 
\begin{equation}
\begin{split}\label{eq: XI and omegaJ}
\langle [X_{J}] , e(V_K) \rangle_{Pet_{n}} 
= \int_{X_{J}}  e(V_K).
\end{split}
\end{equation}
For the case $J=K$, the claim follows from the previous lemma. Let us consider the case $J\ne K$. This condition and $|J|= |K|$ imply that $J\not\supseteq K$. 
We prove that the right hand side of \eqref{eq: XI and omegaJ} is equal to zero due to the vanishing of the Euler class $e(V_K)$ on~$X_{J}$.
Recall from Proposition~\ref{prop: zero locus} that we have the section $\psi_{K}\colon Pet_{n} \rightarrow V_{K}$ satisfying $Z(\psi_K)=\Omega_K$. Thus, the vector bundle $V_{K}|_{X_J}$ restricted to $X_J$ admits a nowhere-zero section given by $\psi_K|_{X_J}$ since we have $X_J\cap\Omega_K=\emptyset$ by $J\not\supseteq K$ and Proposition~\ref{prop: point intersection 2}. Thus the Euler class $e(V_K)$ vanishes on $X_J$, and hence the right hand side of \eqref{eq: XI and omegaJ} is equal to $0$ in this case.
\end{proof}

For $J\subseteq [n-1]$, recall from Definition~\ref{defi:varpi} that
\begin{align*}
\varpi_J 
= \frac{1}{m_J} e(V_J)
= \frac{1}{m_J} \prod_{i\in J}\varpi_i .
\end{align*} 

\begin{theorem}\label{thm: Z-basis of cohomology of Peterson}
For each $J\subseteq [n-1]$, the cohomology class $\varpi_J$ is an element of the integral cohomology $H^{2|J|}(Pet_{n};{\mathbb Z})$, and the set
\begin{align*}
\{ \varpi_J \in H^*(Pet_{n};{\mathbb Z}) \mid J\subseteq [n-1]\}
\end{align*}
is a ${\mathbb Z}$-basis of $H^*(Pet_{n};{\mathbb Z})$.
\end{theorem}

\begin{proof}
Recall from Proposition~\ref{prop: homology basis} that $\{[X_{J}] \mid J\subseteq[n-1]\}$ forms a ${\mathbb Z}$-basis of $H_*(Pet_{n};{\mathbb Z})$. Since the paring between $H_*(Pet_{n};{\mathbb Z})$ and $H^*(Pet_{n};{\mathbb Z})$ is perfect, the previous proposition implies the desired claim.
\end{proof}

\begin{exam}
{\rm 
Let $n=4$ so that $[n-1]=\{1,2,3\}$. The additive basis given in Theorem~\ref{thm: Z-basis of cohomology of Peterson} is
\begin{align*}
H^*(Pet_{n};{\mathbb Z}) &= {\mathbb Z}\varpi_{\emptyset}\oplus({\mathbb Z}\varpi_{\{1\}}\oplus{\mathbb Z}\varpi_{\{2\}}\oplus{\mathbb Z}\varpi_{\{3\}})\\
&\hspace{75pt}\oplus({\mathbb Z}\varpi_{\{1,2\}}\oplus{\mathbb Z}\varpi_{\{1,3\}}\oplus{\mathbb Z}\varpi_{\{2,3\}})\oplus{\mathbb Z}\varpi_{\{1,2,3\}}.
\end{align*}
As we saw in Proposition $\ref{prop: duality}$, this is the dual basis of the basis of the homology group $H_*(Pet_{n};{\mathbb Z})$ given in Example~$\ref{ex: homology basis}$.}
\end{exam}

\subsection{Structure constants and their positivity}\label{subsec: positivity}
By Theorem~\ref{thm: Z-basis of cohomology of Peterson}, we can study the cohomology ring $H^*(Pet_{n};{\mathbb Z})$ in terms of the basis $\{\varpi_J\}_{J\subseteq[n-1]}$. Specifically, we expand the product of two classes $\varpi_J$ and $\varpi_K$ as a linear combination of the basis:
\begin{align} \label{eq:Schubert_calculus_for_varpi_Pet} 
\varpi_J \cdot \varpi_K =\sum_{L \subseteq [n-1]} d_{JK}^L \, \varpi_L, \ \ \ d_{JK}^L \in {\mathbb Z}.
\end{align}
The coefficients $d_{JK}^L$ are called the structure constant for the basis $\{\varpi_J\}_{J\subseteq[n-1]}$. In the following, we explain a geometric interpretation of $d_{JK}^L$, and deduce  their positivity. Note that the degree of $\varpi_L$ in $H^*(Pet_{n};{\mathbb Z})$ is $2|L|$ and that the degree of $\varpi_J \cdot \varpi_K$ in $H^*(Pet_{n};{\mathbb Z})$ is $2(|J|+|K|)$. Thus we may assume that 
\begin{align} \label{eq: degree of piL}
|L|=|J|+|K|
\end{align}
for each summand of \eqref{eq:Schubert_calculus_for_varpi_Pet} since we have $d_{JK}^L =0$ otherwise. Then by Proposition~\ref{prop: duality}, we have
\begin{align}\label{eq: djkl in terms of pairing}
d_{JK}^L 
&= \langle [X_{L}] , \varpi_{J} \cdot \varpi_{K} \rangle_{Pet_{n}} 
= \frac{1}{m_J}
\frac{1}{m_K}\int_{X_L} \left(\prod_{j\in J}\varpi_j\right)\left(\prod_{k\in K}\varpi_k\right) , 
\end{align}
where $m_J$ and $m_K$ are the positive integers defined in \eqref{eq: def of mJ}. Now recall that each $\varpi_i\in H^2(X_L;{\mathbb Z})$ is the Euler class of the line bundle $V_{\{i\}}$ corresponding to the divisor $E_i(=Z(\psi_{\{i\}}))$ on $Pet_{n}$. Hence it follows from \eqref{eq: djkl in terms of pairing} that the structure constant $d_{JK}^L$ computes an intersection number of (possibly duplicate) divisors $E_i\cap X_L$'s on $X_L$ up to a constant multiple given by $\frac{1}{m_J}\frac{1}{m_K}$ (c.f. \cite[Sect.\ 1.1.C]{Lazarsfeld}). This provides a geometric interpretation of $d_{JK}^L$ in \eqref{eq:Schubert_calculus_for_varpi_Pet}, and it leads us to the following instance of positivity.

\begin{prop}\label{prop: positivity}
We have $d_{JK}^L \ge0$ for all $J,K,L\subseteq [n-1]$.
\end{prop}

\begin{proof}
Recall that each line bundle $L_{\varpi_i}$ over $Fl_n$ is nef for $1\le i\le n-1$ (e.g.\ \cite[proof of Proposition 1.4.1]{Bri} or \cite[Lemma 3.5]{ab-fu-ze2}). Hence the restriction of $L_{\varpi_i}$ over $X_L$ is nef as well for $1\le i\le n-1$. Thus the claim $d_{JK}^L \ge 0$ follows from \eqref{eq: djkl in terms of pairing} and the positivity of intersection numbers of nef divisors \cite[Example 1.4.16]{Lazarsfeld}.
\end{proof}
